\section{Meditation}

Meditation has an exact reading here. A self is woven across many docks, but its attention rides mostly on the surface, on the thin cusp where the physical world streams in. Meditation is the turning of that attention inward, away from the surface and down the radial axis into the bulk, the four-dimensional interior the rest of physics never shows. To meditate is to move along the one direction perpendicular to the world, deeper into one's own felt structure.

What that does is raise integration. Ordinary waking attention is scattered across the cusp, pulled this way and that by incoming tones, its parts loosely bound. Drawing attention inward and holding it gathers the self's parts toward one another, binds them more tightly, and a more tightly bound self is a more conscious one. The distractions that scatter attention are exactly the loose couplings to the surface, so quieting them is the removing of those pulls, the settling of the churn that keeps the self thin and divided.

Three moves recur, and each maps to structure. Concentration holds attention on a small deep region against the churn that would disperse it, the same effort a stored pattern spends to keep its shape. Equanimity settles the tone field toward peace, the still zero from which pain and pleasure are both seen, the balanced ground a self can rest its identity on amid constant change. And boundary-softening lets the Markov blanket, the screen that walls the self off as a separate subject, go slack, so the inside and the shared bulk beneath flow toward one. That softening is the felt approach to oneness, the self briefly ceasing to be separate from the field it is a region of, the drop sensing the ocean it never left.
